
\chapter{Hardy \texorpdfstring{$Z$}{Z} 函数、Riemann--Siegel 公式与数值计算}\label{ch:riemann_numerical_approximation}

本章讨论两类计算问题：
（1）用 Hardy $Z$ 函数与 Riemann--Siegel 公式计算非平凡零点虚部（Gram 点、符号变化、割线法精化）；
（2）用有限个零点作 $\pi_0$（及 $\psi_0$）对 $\pi$（及 $\psi$）的数值逼近（实用记 $\pi_K$），并估计逼近误差。

% ============================================================
\section{引言}
% ============================================================

第~\ref{ch:riemann_explicit_formula}~章给出（跳跃点取左右平均，记为 $\psi_0$；记号见《符号与约定》）
\[
\psi_0(x)
=
x
-
\sum_{\rho}\frac{x^{\rho}}{\rho}
-
\log(2\pi)
-
\frac12\log(1-x^{-2}).
\]
数值实现需要：
\begin{enumerate}
  \item 计算非平凡零点 $\rho=\beta+i\gamma$（实践中常在黎曼猜想（RH）工作假设下取 $\beta=1/2$）；
  \item 控制截断到有限个 $\rho$ 时的误差。
\end{enumerate}

第~\ref{sec:hardy_RS}~节建立 Hardy $Z$ 函数与 Riemann--Siegel 公式；
第~\ref{sec:zeros}~节用它们做临界线零点的定位与精化；
第~\ref{sec:explicit}~节给出显式公式的可计算形式、截断参数与程序库入口
（含 $\mathrm{li}$、$\Ei$ 约定）；
第~\ref{sec:numerical}~节给出数值示例。

% ============================================================
\section{Hardy \texorpdfstring{$Z$}{Z} 函数与 Riemann--Siegel 公式}\label{sec:hardy_RS}
% ============================================================

\paragraph{知识谱系定位}
在全书链条
$\Gamma$--函数方程--临界线--零点--显式公式
中，本节对象用于\textbf{衔接}上述环节，本身并非新的 $L$-函数：
\begin{itemize}
  \item $\vartheta(t)$：第~\ref{ch:gamma_function}~、\ref{ch:riemann_zeta_continuation}~章中
        完备化因子 $\pi^{-s/2}\Gamma(s/2)$ 在临界线 $s=\frac12+it$ 上的\textbf{相位}；
  \item $Z(t)$：临界线上的\textbf{实振幅}（临界线实化），使求零化为实函数 $Z$ 的变号；
  \item Riemann--Siegel 公式：近似函数方程在临界线上的\textbf{高效求值}，服务第~\ref{sec:zeros}~节，
        而非替代解析延拓。
\end{itemize}
二者亦为第~\ref{ch:visualization}~章螺旋几何的同一分解。
下文给出 $Z$ 实值性的完整短证，以及 Riemann--Siegel 主项的来源提要；
精细余项的一致估计外引经典文献（见各小节末）。

\subsection{Hardy \texorpdfstring{$Z$}{Z} 函数}

本书前文章节的主对象是 $\zeta(s)$；临界线上 $\zeta(\frac12+it)$ 一般仍为\textbf{复数}，
其实部与虚部同时振荡，直接搜索零点相当于同时解两个实方程。
本小节引入的 Hardy $Z$ 函数并非新的 $L$-函数，而是对 $\zeta(\frac12+it)$ 作\textbf{相位归一}
后得到的\textbf{实值}辅助函数，专用于临界线零点的数值计算与几何解读
（第~\ref{ch:visualization}~章的螺旋轨迹即同一分解的几何示意）。

\begin{definition}{Riemann--Siegel 相位 $\vartheta(t)$}{}
\[
\vartheta(t)
=
\Imag\!\left[\log\Gamma\!\left(\frac14+\frac{it}{2}\right)\right]
-
\frac{t}{2}\log\pi.
\]
\end{definition}

此处 $\log\Gamma(\cdot)$ 取沿正实轴连续的主值分支；对 $t<0$ 有
$\vartheta(-t)=-\vartheta(t)$（奇函数）。通常只需考虑 $t>0$。


\paragraph{记号区分}
本书中三种易混的“theta”严格分工如下（全文统一）：
\begin{itemize}
  \item $\theta(x)$：\textbf{第一 Chebyshev 函数}（第~\ref{ch:elementary_number_theoretic_functions}~章），
        自变量为 $x>0$ 的计数变量；
  \item $\vartheta(t)$：\textbf{Riemann--Siegel 相位}（本定义），
        自变量为临界线虚部 $t\in\mathbb{R}$；
  \item Jacobi $\vartheta$（常写 $\vartheta(x)$ 或 $\vartheta(t)$，$t>0$）：
        第~\ref{ch:riemann_zeta_continuation}~、\ref{ch:integral_transforms}~章模反演用的
        $\sum_n e^{-\pi n^2 t}$，与本相位\textbf{不是同一对象}。
\end{itemize}


$\vartheta(t)$ 的渐近展开（由 Stirling 公式导出，第~\ref{ch:gamma_function}~章）：
\begin{equation}
  \vartheta(t) = \frac{t}{2}\log\frac{t}{2\pi} - \frac{t}{2} - \frac{\pi}{8}
              + \frac{1}{48t} + \frac{7}{5760 t^3} + O(t^{-5}).
  \label{eq:theta_asymp}
\end{equation}

\begin{definition}{Hardy $Z$ 函数}{}
\[
Z(t) = e^{i\vartheta(t)}\,\zeta\!\left(\frac12+it\right).
\]
\end{definition}

\paragraph{$Z(t)$ 的实值性（正文短证）}
记完备化因子（第~\ref{ch:xi_function}~章记号）
\[
\xi^*(s)=\pi^{-s/2}\,\Gamma\!\left(\frac{s}{2}\right)\,\zeta(s)
\]
（与 $\xi(s)$ 仅差整函数因子 $\frac12 s(s-1)$，不影响此处相位关系）。
由第~\ref{ch:riemann_zeta_continuation}~章函数方程~\eqref{eq:asymmetric_fe}--\eqref{eq:asymmetric_fe_alt}
等价地有 $\xi^*(s)=\xi^*(1-s)$。取 $s=\frac12+it$（$t\in\mathbb{R}$），得
\begin{equation}
  \pi^{-\frac14-\frac{it}{2}}\,
  \Gamma\!\left(\frac14+\frac{it}{2}\right)\,
  \zeta\!\left(\frac12+it\right)
  =
  \pi^{-\frac14+\frac{it}{2}}\,
  \Gamma\!\left(\frac14-\frac{it}{2}\right)\,
  \zeta\!\left(\frac12-it\right).
  \label{eq:xi_star_crit}
\end{equation}
整理得
\begin{equation}
  \zeta\!\left(\frac12+it\right)
  =
  \pi^{it}\,
  \frac{\Gamma\!\left(\frac14-\frac{it}{2}\right)}{\Gamma\!\left(\frac14+\frac{it}{2}\right)}\,
  \zeta\!\left(\frac12-it\right).
  \label{eq:zeta_crit_fe}
\end{equation}
因 $\Gamma(\overline{z})=\overline{\Gamma(z)}$，有
$\Gamma(\frac14-\frac{it}{2})=\overline{\Gamma(\frac14+\frac{it}{2})}$，故
\[
\frac{\Gamma\!\left(\frac14-\frac{it}{2}\right)}{\Gamma\!\left(\frac14+\frac{it}{2}\right)}
=
\exp\!\Bigl(-2i\,\Imag\log\Gamma\!\left(\frac14+\frac{it}{2}\right)\Bigr),
\qquad
\pi^{it}=e^{it\log\pi}.
\]
合并指数并与 $\vartheta$ 的定义对照，即得临界线上的相位形式
\begin{equation}
  \zeta\!\left(\frac12+it\right)
  =
  e^{-2i\vartheta(t)}\,
  \zeta\!\left(\frac12-it\right).
  \label{eq:zeta_phase_fe}
\end{equation}
又因 $\zeta$ 的 Dirichlet 系数（从而 Euler 乘积与解析延拓）在实轴上取实值，
Schwarz 反射给出 $\zeta(\overline{s})=\overline{\zeta(s)}$，故
$\zeta(\frac12-it)=\overline{\zeta(\frac12+it)}$。代入 \eqref{eq:zeta_phase_fe} 得
\begin{equation}
  \zeta\!\left(\frac12+it\right)
  =
  e^{-2i\vartheta(t)}\,
  \overline{\zeta\!\left(\frac12+it\right)}.
  \label{eq:zeta_conj_crit}
\end{equation}
现在计算 $Z$ 的复共轭：
\begin{align*}
  \overline{Z(t)}
  &=
  e^{-i\vartheta(t)}\,
  \overline{\zeta\!\left(\frac12+it\right)}
  =
  e^{-i\vartheta(t)}\,
  e^{2i\vartheta(t)}\,
  \zeta\!\left(\frac12+it\right)
  =
  e^{i\vartheta(t)}\,
  \zeta\!\left(\frac12+it\right)
  =
  Z(t),
\end{align*}
其中第二步用了 \eqref{eq:zeta_conj_crit}。因此 $Z(t)\in\mathbb{R}$。
又 $|e^{i\vartheta(t)}|=1$，故
\[
  |Z(t)|=\bigl|\zeta\!\left(\frac12+it\right)\bigr|.
\]
（上述推导与 Titchmarsh (1986) §4.17 的标准叙述一致，可供对照；此处已在本书框架内写完。）

\begin{corollary}{}{}
$\zeta(\frac12+it_0)=0$ 当且仅当 $Z(t_0)=0$。
\end{corollary}
因此，寻找临界线上零点等价于寻找实函数 $Z(t)$ 的零点。

\paragraph{$Z$ 与 $\zeta$ 的对照（必要说明）}
由定义立刻反解出 \textbf{Hardy 极坐标形式}
\begin{equation}
  \zeta\!\left(\frac12+it\right)
  =
  Z(t)\,e^{-i\vartheta(t)}.
  \label{eq:hardy_polar}
\end{equation}
其几何读法是：
\begin{itemize}
  \item $e^{i\vartheta(t)}$（或 \eqref{eq:hardy_polar} 中的 $e^{-i\vartheta(t)}$）
        是模长为 $1$ 的相位因子——在复平面上作纯旋转；
  \item $Z(t)\in\mathbb{R}$ 是\textbf{带符号的实振幅}（可正可负），
        几何半径为 $|Z(t)|=\bigl|\zeta(\frac12+it)\bigr|$，
        负号等价于再转 $\pi$；
  \item 故 $Z(t)\neq\zeta(\frac12+it)$：后者一般是复数，前者是把后者
        旋到实轴上之后的实坐标。
\end{itemize}
换言之，$Z$ \textbf{不改写} $\zeta$ 的定义，而只是把问题限制在临界线
$s=\frac12+it$ 上：将「寻找复值函数 $\zeta(\frac12+it)$ 的零点」
化为「寻找实函数 $Z(t)$ 的变号点」。
第~\ref{sec:RS}~节的 Riemann--Siegel 公式
则是对 $Z(t)$ 的高效近似，而非对 $\zeta$ 解析延拓的替代。
表~\ref{tab:Z_vs_zeta} 汇总对照。

\begin{table}[htbp]
\centering
\caption{$Z(t)$ 与 $\zeta(\frac12+it)$ 对照}
\label{tab:Z_vs_zeta}
\begin{tabular}{@{}lll@{}}
\toprule
对象 & 类型 & 本书中的角色 \\
\midrule
$\zeta(s)$ & 亚纯（主对象） & 前文章节主线：定义、函数方程、零点、显式公式 \\
$\zeta(\frac12+it)$ & 复值 & 临界线上的限制；螺旋轨迹本身 \\
$\vartheta(t)$ & 实值相位 & 由 $\Gamma$ 决定的旋转角（非 Jacobi $\vartheta$） \\
$Z(t)$ & 实值 & 振幅（可正负）；求零点；Hardy 极坐标形式 \\
$|Z(t)|$ & 非负实数 & $\bigl|\zeta(\frac12+it)\bigr|$，螺圈半径 \\
\bottomrule
\end{tabular}
\end{table}

\subsection{Riemann--Siegel 公式}\label{sec:RS}

直接计算 $Z(t)$ 需要求 $\zeta(\frac12+it)$；若对 $\Real s>1$ 的 Dirichlet 级数作解析延拓后硬算，
项数随 $t$ 增大而剧增。Riemann--Siegel 公式把\textbf{近似函数方程}限制到临界线，
给出 $O(\sqrt{t})$ 复杂度的求值，是第~\ref{sec:zeros}~节大规模数零点的引擎。

\paragraph{主项来源提要（正文短述）}
由非对称函数方程~\eqref{eq:asymmetric_fe_alt}，在远离极点处可将 $\zeta(s)$ 写成
「长度约 $\sqrt{|t|/(2\pi)}$ 的 Dirichlet 截断」与「反射项 $\chi(s)$ 乘以 $\zeta(1-s)$ 的同类截断」
之和，再加上可控余项；这就是\textbf{近似函数方程}（approximate functional equation）的轮廓。
取 $s=\frac12+it$，并令截断参数
\[
  N_{\mathrm{RS}}=\Bigl\lfloor\sqrt{\frac{t}{2\pi}}\Bigr\rfloor,
\]
则两侧截断长度相同。将 \eqref{eq:zeta_phase_fe} 的相位因子并入 $e^{i\vartheta(t)}$ 后，
两项合并为同一实余弦和：主贡献正是
\[
  2\sum_{n=1}^{N_{\mathrm{RS}}}
  \frac{\cos\!\bigl(\vartheta(t)-t\log n\bigr)}{\sqrt{n}},
\]
余项在标准估计下为 $O(t^{-1/4})$。由此得到下述可计算形式。

\begin{theorem}{Riemann--Siegel 公式}{}
令 $N_{\mathrm{RS}} = \lfloor\sqrt{t/(2\pi)}\rfloor$。则
\begin{equation}
  Z(t) = 2\sum_{n=1}^{N_{\mathrm{RS}}} \frac{\cos\!\bigl(\vartheta(t)-t\log n\bigr)}{\sqrt{n}} + R(t),
  \label{eq:RS}
\end{equation}
其中余项 $R(t) = O(t^{-1/4})$。更精确地，设 $p = \{\sqrt{t/(2\pi)}\}$（分数部分），
首项修正常取
\[
R(t) = (-1)^{N_{\mathrm{RS}}-1}\left(\frac{2\pi}{t}\right)^{1/4}
       \left[ \frac{\cos\bigl(2\pi(p^2-p-1/16)\bigr)}{\cos(2\pi p)} + O(t^{-1/2}) \right].
\]
\end{theorem}

\paragraph{证明与文献分工}
\begin{itemize}
  \item \textbf{本书承担}：上段主项来源（近似函数方程 $\to$ 临界线余弦和）、$N_{\mathrm{RS}}$ 的尺度，
        以及 \eqref{eq:RS} 作为零点算法的输入；
  \item \textbf{完整渐近与一致余项}：见 Edwards, \emph{Riemann's Zeta Function}
        （相关章节）或 Titchmarsh (1986) Chapter~4；
  \item \textbf{首项修正在退化点的处理}：当 $\cos(2\pi p)=0$（即 $p=\tfrac14$ 或 $\tfrac34$）时
        上列 $R(t)$ 表达式分母为零，需改用更高阶 Riemann--Siegel 展开，
        见 Turing (1953) 或 Berry \& Keating (1992)——二者针对\textbf{余项精化}，
        而非替代主公式的出处。
\end{itemize}

该公式的计算量为 $O(\sqrt{t})$，远优于对长度 $\sim t$ 的直接截断，
使得在个人计算机上计算 $t\sim 10^6$ 量级的零点成为可能。
（注：上式中 $N_{\mathrm{RS}}$ 为 Riemann--Siegel 截断参数，与本章第~\ref{sec:explicit}~节的 Möbius 截断参数 $N=\lfloor \log_2 x \rfloor + 1$ 含义不同，请注意区分。）

% ============================================================
\section{非平凡零点的数值计算}\label{sec:zeros}
% ============================================================

在第~\ref{sec:hardy_RS}~节的 $Z(t)$ 与 Riemann--Siegel 公式之下，
本节汇总零点计算所需的记号与计数函数，并给出 Gram 点定位与割线法精化。
（对数积分与指数积分仅用于将零点代入显式公式，见第~\ref{sec:explicit}~节，此处不引入；
程序库、预计算数据与代表性验证结果见该节末。）

\subsection{基本符号与零点计数}\label{sec:prelim}

$\zeta(s)$ 的非平凡零点记为
$\rho = \beta + i\gamma$，其中 $0<\beta<1$，$\gamma\in\R$。
由第~\ref{ch:zero_distribution}~章，非平凡零点关于实轴和临界线 $\Real(s)=\frac12$ 均对称分布。
在 RH（未证明，但本章数值计算中作为工作假设）下，
$\rho_k = \frac12 + i\gamma_k$，$\gamma_k > 0$。

\begin{definition}{零点计数函数}{}
$N(T) = \#\{\rho: \zeta(\rho)=0,\ 0<\Imag(\rho)\le T\}$（计重数）。
\end{definition}

由第~\ref{ch:zero_distribution}~章 von Mangoldt 公式：
\begin{equation}
  N(T) = \frac{T}{2\pi}\log\frac{T}{2\pi e} + O(\log T).
  \label{eq:NT_chap15}
\end{equation}

表~\ref{tab:zeros} 列出前十个非平凡零点虚部（精确至六位小数），作为后文算法与数值实验的对照基准。

\begin{table}[htbp]
\centering
\caption{$\zeta(s)$ 前十个非平凡零点虚部 $\gamma_k$}
\label{tab:zeros}
\begin{tabular}{@{}crcr@{}}
\toprule
$k$ & $\gamma_k$ & $k$ & $\gamma_k$ \\
\midrule
1 & 14.134725 & 6 & 37.586178 \\
2 & 21.022040 & 7 & 40.918720 \\
3 & 25.010858 & 8 & 43.327073 \\
4 & 30.424876 & 9 & 48.005151 \\
5 & 32.935062 & 10 & 49.773832 \\
\bottomrule
\end{tabular}
\end{table}

\subsection{Gram 点与零点定位}

\begin{definition}{Gram 点}{}
满足 $\vartheta(g_n) = n\pi$ 的实数 $g_n$ 称为第 $n$ 个 Gram 点。
\end{definition}

由渐近展开 \eqref{eq:theta_asymp} 知，当 $t$ 充分大时（约 $t>2\pi e\approx17.08$）$\vartheta(t)$ 严格单调递增，故对足够大的 $n$，Gram 点存在唯一；对有限个小 $n$（如 $n=-1,0$）则通过直接计算确认。表~\ref{tab:gram} 列出前几个 Gram 点。

\begin{table}[htbp]
\centering
\caption{前几个 Gram 点}
\label{tab:gram}
\begin{tabular}{@{}ccc@{}}
\toprule
$n$ & $g_n$ & $(-1)^n Z(g_n)$ 符号 \\
\midrule
$-1$ & 9.6669 & $+$ \\
0 & 17.8456 & $+$ \\
1 & 23.1703 & $+$ \\
2 & 27.6702 & $+$ \\
3 & 31.7198 & $+$ \\
4 & 35.4671 & $+$ \\
\bottomrule
\end{tabular}
\end{table}

\begin{definition}{Gram 法则}{}
若 $(-1)^n Z(g_n) > 0$，称 $g_n$ 为\textbf{好 Gram 点}。
约 $72.7\%$ 的 Gram 点满足此性质。
\end{definition}

零点搜索算法如下：

\begin{alg}{Gram 点定位}{}
\label{alg:gram}
给定 $T>0$：
\begin{enumerate}
  \item 计算 Gram 点 $g_0,g_1,\ldots$ 直至超过 $T$。
  \item 对每个区间 $[g_{n-1},g_n]$，用 \eqref{eq:RS} 计算 $Z(g_{n-1})$ 与 $Z(g_n)$。
  \item 若 $Z(g_{n-1})=0$ 或 $Z(g_n)=0$，则对应端点本身即为零点，直接精化（此情况极为罕见，实践中通过微扰处理）。
  \item 若 $Z(g_{n-1})$ 与 $Z(g_n)$ 符号相反，则该区间存在奇数个零点；否则需细分搜索。
  \item 用二分法将零点区间缩小至精度 $\varepsilon_0$。
\end{enumerate}
\end{alg}

\subsection{割线法精化}

由于 $Z'(t)$ 计算代价较高，实践中常用\textbf{割线法}（Secant Method）代替 Newton 法进行精化。割线法用差商近似导数，收敛阶约为 $1.618$（低于 Newton 法的二阶）：

得到初始近似 $t_0$ 后，迭代格式为：
\begin{equation}
  t_{k+1} = t_k - Z(t_k)\,\frac{t_k-t_{k-1}}{Z(t_k)-Z(t_{k-1})}.
  \label{eq:secant}
\end{equation}

\begin{example}{}{}
在 $(9.667,17.846)$ 内搜索第一个零点：二分法得 $t_0\approx14.13$，
经五步割线迭代收敛至
\[
\gamma_1 \approx 14.134725141734693\ldots
\]
误差 $<10^{-15}$。
\end{example}

% ============================================================
\section{显式公式的数值实现}\label{sec:explicit}
% ============================================================

将第~\ref{sec:zeros}~节求得的零点代入显式公式以逼近 $\pi_0$ 或 $\psi_0$
（及实用截断 $\pi_K$）时，需要对数积分与指数积分；
二者\textbf{不参与}临界线零点本身的定位与计数，
故在本节而非第~\ref{sec:zeros}~节引入。

\subsection{对数积分与指数积分}

定义、实轴恒等式、复 $\li(x^{\rho})$ 与全书分工见第~\ref{ch:logint_Ei}~章
（第~\ref{sec:Li}~、\ref{sec:Ei}~、\ref{sec:Li-li-convention}~节）。
本节只调用
\[
  \li(x)=\Ei(\log x),\qquad
  \li(x^{\rho})\longleftrightarrow\Ei(\rho\log x)
\]
作截断计算；\textbf{显式公式叙述用 $\li$}，\textbf{$\pi\sim\Li$ 比较用 $\Li$}。

\subsection{显式公式的 \texorpdfstring{$\Ei$}{Ei} 形式}

第~\ref{ch:riemann_explicit_formula}~章中的 Riemann--von Mangoldt 显式公式~\eqref{eq:explicit_psi} 给出
（跳跃点为 $\psi_0$；记号见《符号与约定》）：
\[
\psi_0(x) = x - \sum_{\rho} \frac{x^{\rho}}{\rho} - \log(2\pi) - \frac12\log(1-x^{-2}).
\]

通过分部积分（见第~\ref{ch:riemann_explicit_formula}~章第~\ref{sec:psi-to-pi}~节相关推导），可得 $J$ 的显式公式
（主项用 $\li$，非偏移 $\Li$）：
\begin{equation}
  J_0(x) = \li(x) - \sum_{\rho} \li(x^{\rho}) - \log 2 + \int_x^\infty \frac{dt}{t(t^2-1)\log t}.
  \label{eq:explicit_J_chap15}
\end{equation}

利用 $\li(x) = \Ei(\log x)$，可将上式改写为纯 $\Ei$ 形式。
常数项 $-\log 2$ 经 Möbius 反演形式上对应 $-\log 2\sum_n\mu(n)/n$，
在 $s=1$ 处与 $1/\zeta(1)=0$ 相关，但该级数仅条件收敛/需正则化；
有限截断 $N_\mu=\lfloor\log_2 x\rfloor+1$ \textbf{并不会}精确抵消该常数。
对 $x\in[10,40]$ 量级的示意作图，常数误差常小于绘图精度，故实用公式中常省略；
精确比较时应保留或估计残差 $\sum_{n\le N_\mu}\mu(n)/n$。得到数值计算实用的形式：

\begin{equation}\label{eq:pi_explicit_full}
\boxed{\begin{aligned}
\pi_0(x) &\approx \sum_{n=1}^{N_{\mu}} \frac{\mu(n)}{n} \Ei\!\left(\frac{\log x}{n}\right) \\
&\quad - 2\sum_{k=1}^{K} \Re \left[ \sum_{n=1}^{N_{\mu}} \frac{\mu(n)}{n} \Ei\!\left(\frac{\rho_k \log x}{n}\right) \right] \\
&\quad - \sum_{m=1}^{M} \sum_{n=1}^{N_{\mu}} \frac{\mu(n)}{n} \Ei\!\left(\frac{-2m \log x}{n}\right).
\end{aligned}}
\end{equation}

其中：
\begin{itemize}
  \item $\mu(n)$ 为 Möbius 函数；
  \item $\rho_k = \frac12 + i\gamma_k$ 为前 $K$ 个非平凡零点；
  \item 第三项对应 $\zeta(s)$ 在 $s=-2m$ 处的平凡零点贡献；
  \item $N_{\mu}$、$M$、$K$ 为截断参数。
\end{itemize}

\subsection{公式的结构解析}

\begin{itemize}
  \item \textbf{第一项（光滑主项）}：$\sum_{n=1}^{N_{\mu}} \frac{\mu(n)}{n} \Ei(\frac{\log x}{n})$，
        即 Riemann $R(x)$ 函数的截断，提供 $\pi(x)$ 的主要近似。
  \item \textbf{第二项（非平凡零点修正）}：$-2\sum_{k=1}^{K}\Re[\cdots]$，
        每个零点贡献一个频率为 $\gamma_k$ 的振荡。由第~\ref{ch:riemann_explicit_formula}~章的振幅分析：
        对 $\psi_0$，共轭对合并后振幅为 $2x^{1/2}/|\rho_k|$（若 RH 成立）；
        对 $\pi_0$（或 $J_0$），振幅约为 $2x^{1/2}/(|\rho_k|\log x)$，
        这来自 $\Li$ 的渐近展开 $\Li(x^\rho)\sim x^\rho/(\rho\log x)$（$|\rho|\to\infty$）。
  \item \textbf{第三项（平凡零点修正）}：对应负偶整数零点，随 $m$ 增大迅速衰减，
        通常取 $M=5\sim10$ 即可。
\end{itemize}

\subsection{参数选择与误差分析}

由第~\ref{ch:zero_distribution}~章 von Mangoldt 公式，高度不超过 $T$ 的零点个数 $N(T) \sim \frac{T}{2\pi}\log T$。

由第~\ref{ch:riemann_explicit_formula}~章 $\psi_0$ 显式公式的截断误差分析（RH 成立时），
更稳妥的是用截断高度 $T$ 陈述尾项。设取 $|\Im\rho|\le T$ 的对称截断，则典型估计为
\[
\bigl|\psi_0(x)-\psi_T(x)\bigr|
=
O\!\left(
\frac{x^{1/2}\log^2(xT)}{T}
\right)
\qquad(x\ge2,\;T\ge2),
\]
（常数依赖所引用的零点密度/$\zeta'/\zeta$ 界；此处作\textbf{阶估计/启发式}，不写成定理级 $\lesssim$；
$\psi_T$ 表截断到高度 $T$ 的部分和）。
由 $N(T)\sim\frac{T}{2\pi}\log\frac{T}{2\pi}$，若 $K=N(T)$ 为所用零点个数，则
$T\sim 2\pi K/\log K$，代入得
\[
\bigl|\psi_0(x)-\psi_K(x)\bigr|
=
O\!\left(
\frac{x^{1/2}\log K\cdot\log^2(xK)}{K}
\right).
\]
对 $\pi$ 或 $J$ 尚需经 Abel/Möbius 变换，主振荡振幅多一个 $1/\log x$ 因子。
若要求 $|\psi_0-\psi_K|<\varepsilon$，则大致需 $T\gg x^{1/2}\log^2(xT)/\varepsilon$，
再换算为 $K=N(T)$；\textbf{不可}将 $K\gtrsim C_1 x^{1/2}\log K/\varepsilon$ 当作精确充分条件。

Möbius 截断参数取 $N_{\mu} = \lfloor \log_2 x \rfloor + 1$ 即可保证 $x^{1/N_{\mu}} < 2$，
此时 $J(x^{1/N_{\mu}})=0$，更高阶贡献可忽略。

平凡零点截断 $M=10$ 已足够，因为 $\Ei(-2m\log x/n)$ 随 $m$ 指数衰减。

\subsection{数值计算的程序库与平台}\label{sec:libraries}

前文给出了求零点与 $\pi_0$ 逼近的\textbf{数学流程}
（Hardy $Z$、Riemann--Siegel、Gram、割线；Möbius--$\Ei$ 与参数 $N_{\mu},K,M$）。
实现时还须选定求值引擎：直接编码 RS 余弦和可行，但大规模或需严格误差时，
通常调用专用库或读取预计算零点表。
本节按\textbf{任务分层}介绍常用工具，并给出与本章算法的对应；
命令行与代码片段仅作入口示意，完整手册见各项目文档。
安装方式随平台而异（包管理器、源码编译或语言包索引），此处不展开。

\subsubsection{任务分层与选用}

本章相关计算可粗分为四类，对应不同工具：

\begin{center}
\renewcommand{\arraystretch}{1.25}
\begin{tabular}{@{}p{3.2cm}p{4.0cm}p{5.8cm}@{}}
\toprule
\textbf{任务} & \textbf{典型需求} & \textbf{常用工具} \\
\midrule
临界线求值 / 找零点
  & $Z(t)$ 或 $\zeta(\tfrac12+it)$ 快速求值、变号定位
  & lcalc；自写 RS（第~\ref{sec:RS}~节）；PARI/GP \\
高精度浮点
  & 任意位数 $+\times\div$、初等函数
  & MPFR（C）；mpmath（Python） \\
严格误差 / 区间证明
  & 包围区间保证含真值；有限高度 RH 验证
  & Arb（现并入 FLINT）；arb 系 Python 绑定 \\
批量数据与作图
  & 已知 $\gamma_k$ 表、LMFDB、画 $\pi_K$
  & Odlyzko 表、LMFDB；mpmath；通用绘图库 \\
\bottomrule
\end{tabular}
\end{center}

\noindent
原则简述：
\begin{itemize}
  \item \textbf{复现表~\ref{tab:zeros} 或小 $K$ 实验}：mpmath 的 \texttt{zetazero} 或预计算表通常足够；
  \item \textbf{按第~\ref{sec:zeros}~节自行实现 Gram--割线}：需稳定的 $Z(t)$ 求值——
        可调用 lcalc / PARI，或按 \eqref{eq:RS} 自写主项并控制余项；
  \item \textbf{声称「某高度内全部零点在临界线上」且要可检验的包围}：
        须区间算术（Arb/FLINT 一脉），而非仅双精度浮点比较 $|\zeta|$ 是否「很小」；
  \item \textbf{第~\ref{sec:numerical}~节波形图}：零点虚部可来自表或库，
        主计算量在 $\Ei(\rho\log x/n)$ 的截断求和。
\end{itemize}

\begin{table}[htbp]
\centering
\caption{与本章相关的主流库（按角色，非完备目录）}
\label{tab:libraries}
\small
\begin{tabular}{@{}p{2.3cm}p{1.6cm}p{9.2cm}@{}}
\toprule
名称 & 形态 & 与本章的关系 \\
\midrule
lcalc
  & C++ / 命令行
  & Rubinstein 的 $L$-函数包；内置 $\zeta$ 与零点搜索，适合批量 $\gamma_k$ \\
PARI/GP
  & C / 交互 \texttt{gp}
  & 通用数论系统；\texttt{lfunzeros} 等可算 $\zeta$ 与 Dirichlet $L$ 的临界零点 \\
MPFR
  & C 库
  & 任意精度浮点（IEEE 风格舍入）；常作更高层库的底座 \\
FLINT / Arb
  & C 库
  & FLINT：整数与多项式等；\textbf{Arb}：球算术（ball arithmetic）。
    FLINT~3 起 Arb 并入 FLINT。用于严格包围与 Platt--Trudgian 类验证的底层 \\
mpmath
  & Python
  & 纯 Python 任意精度；\texttt{zetazero(k)} 取第 $k$ 个零点，便于脚本与作图 \\
SciPy
  & Python
  & \texttt{scipy.special.zeta} 可估复点上的 $\zeta$；\textbf{不}提供系统零点表或严格包围 \\
\bottomrule
\end{tabular}
\end{table}

\subsubsection{专用与交互式系统}

\paragraph{lcalc}
Michael Rubinstein 的 \textbf{lcalc} 面向 $L$-函数的数值计算：特殊值、临界线求值与零点定位。
对本书而言，其角色是把第~\ref{sec:hardy_RS}--\ref{sec:zeros}~节的流程封装为可调用程序——
内部使用针对 $L$-函数优化的求值（含与 RS 同类的思想），用户不必手写余弦和即可得到 $\gamma_k$。
典型命令行用法：
\begin{verbatim}
lcalc --zeros 1000
lcalc --zeros 5 --N 5
lcalc --value --x=0.5 --y=14.134725141738
\end{verbatim}
分别表示：前 $1000$ 个零点；自第 $6$ 个起取 $5$ 个；以及求 $\zeta(1/2+14.1347\ldots i)$。
具体选项以当前版本说明为准。

\paragraph{PARI/GP}
\textbf{PARI/GP} 是解析数论中广泛使用的交互式与可编程系统（\texttt{gp} 为交互前端）。
除 $\zeta$ 外，可统一处理 Dirichlet $L$-函数等对象，便于对照「$\zeta$ 只是特例」。
例如计算 $L(s,\chi_{-4})$（或文档约定的等价写法）在高度 $20$ 以下的临界零点：
\begin{verbatim}
gp > lfunzeros(x^2+1, 20)
\end{verbatim}
对 $\zeta$ 本身亦有 \texttt{zeta}、\texttt{lfun} 族命令。
PARI 侧重通用数论与灵活实验；超大规模零点或可机器检验的区间证明，
通常仍依赖专用实现或 Arb 类库。

\subsubsection{高精度与严格算术}

\paragraph{MPFR}
\textbf{MPFR} 提供正确舍入的任意精度浮点运算，是许多数论程序的底层依赖。
单独使用 MPFR 可实现高精度 $\Gamma$、$\log$ 等，从而自建 $Z(t)$；
工作量与数值稳定性需自行负责。本章实验若只需「足够多位小数的 $\gamma_k$」，
往往不必直接写 MPFR，而通过 mpmath 或更高层库间接使用同类能力。

\paragraph{Arb 与 FLINT}
\textbf{Arb}（Fredrik Johansson）实现\textbf{球算术}：每个量表示为中心 $+$ 半径，
运算时自动传播误差上界，使「结果区间必含真值」在机器上可检验
（在实现与舍入模型正确的前提下）。
这与「双精度算出 $|\zeta(1/2+it)|<10^{-12}$ 因而当作零点」有本质区别：
后者是启发式；前者可支撑有限高度上的\textbf{严格}结论。
第~\ref{ch:zero_distribution}~章与下表中 Platt--Trudgian 的验证，即依赖区间/围道与此类算术。

自 \textbf{FLINT~3} 起，Arb 的源码与接口并入 \textbf{FLINT}（Fast Library for Number Theory）；
新代码以 \texttt{flint/arb.h}、\texttt{flint/acb.h} 等头文件使用球算术，
无需再单独安装历史独立包 \texttt{libarb}（旧文档仍可能分列二者）。
Python 绑定方面，\textbf{python-flint} 等可在脚本中调用 \texttt{arb}/\texttt{acb} 类型；
具体安装见各发行版或语言包索引。

\subsubsection{Python 脚本与作图}

\paragraph{mpmath}
\textbf{mpmath} 在 Python 中提供任意精度浮点与常见特殊函数，并内置 Riemann $\zeta$ 的零点数据接口。
适合：核对表~\ref{tab:zeros}、生成第~\ref{sec:numerical}~节所需的有限个 $\gamma_k$、
以及与 \texttt{matplotlib} 等联用作图。示例：
\begin{verbatim}
from mpmath import zetazero
gamma_1 = zetazero(1).imag   # 14.134725141734693...
zeros = [zetazero(k) for k in range(1, 11)]
\end{verbatim}
默认精度可调（\texttt{mp.dps}）；增大精度会增加耗时。
mpmath 的零点来自内置/预计算数据与求值，\textbf{不等价}于用户完成了一次 Platt--Trudgian 式严格验证。

\paragraph{SciPy 等}
SciPy 的 \texttt{scipy.special.zeta} 可对复自变量求 $\zeta(s)$ 的浮点近似，
用于抽查「表中 $\rho$ 附近 $|\zeta(\rho)|$ 是否接近 $0$」。
它\textbf{不是}零点搜索引擎，也不给出误差的区间证明。
大规模 $L$-函数工作仍应回到 lcalc、PARI 或专用实现。

\paragraph{预计算数据与数据库}
当 $K$ 达数百、数千时，重复现场求零点往往不如使用权威预计算表：
\begin{itemize}
  \item \textbf{Odlyzko} 等发布的零点表（及封装包，如 Python 生态中的相关 data 包）：
        直接提供 $\gamma_k$，供 $\pi_0$ 逼近与统计检验；
  \item \textbf{LMFDB}（L-functions and Modular Forms Database，\texttt{https://www.lmfdb.org}）：
        在线查询 $L$-函数的零点与元数据，并标注验证状态，便于对照与引用。
\end{itemize}
第~\ref{sec:numerical}~节四联图所依赖的 $\gamma_k$，即属「先取零点、再代入 \eqref{eq:pi_explicit_full}」这一路径。

\subsubsection{与本章算法的对应}

\begin{center}
\renewcommand{\arraystretch}{1.3}
\begin{tabular}{@{}p{4.6cm}p{8.4cm}@{}}
\toprule
\textbf{本章步骤} & \textbf{实现时常见做法} \\
\midrule
求 $\vartheta(t)$、$Z(t)$（\eqref{eq:theta_asymp}、\eqref{eq:RS}）
  & 自写 RS 主项；或调用 lcalc / PARI 的临界线求值 \\
Gram 点与割线（第~\ref{sec:zeros}~节）
  & 自写外层循环；内层 $Z$ 求值外包给库 \\
$\gamma_k$ 列表（$K\le 500$ 等）
  & \texttt{zetazero} / 预计算表 / lcalc --zeros \\
$\Ei(\rho\log x/n)$ 与 $\pi_K$
  & mpmath 或任意精度 $\Ei$；按 \eqref{eq:pi_explicit_full} 截断求和 \\
有限高度「全部在临界线上」
  & Arb/FLINT 类区间方法 + 文献级实现；非本章实验范围 \\
\bottomrule
\end{tabular}
\end{center}

\noindent
因此：\textbf{17.2--17.4 前半是数学规格}；\textbf{本小节是实现图层}；
\textbf{第~\ref{sec:numerical}~节是按规格跑通的数值展示}。
三者不可互相替代——库不能代替 RS 公式的推导，图也不能代替有限全称验证的逻辑地位。

\subsubsection{代表性计算结果}

第~\ref{ch:zero_distribution}~章已强调：RH 是全称命题，一反例即否证，
有限高度上的验证不能代替证明，却为搜寻反例与建立有限全称结论提供了理论必要性。
下表所列工作，正是在该逻辑下用不断增强的机器与算法推进「可达到高度」；
「严格验证」与「数值计算」性质不同（见表示后说明）。

\begin{table}[htbp]
\centering
\caption{零点计算的若干结果}
\label{tab:milestones}
\begin{tabular}{@{}cccl@{}}
\toprule
年份 & 研究者 & 零点个数 & 性质 \\
\midrule
1936 & Titchmarsh & $1.0\times 10^3$ & 计算 \\
1956 & Turing & $1.1\times 10^3$ & 计算 \\
1979 & Brent & $8.1\times 10^7$ & 计算 \\
1980s & Odlyzko & $1.5\times 10^9$ & 计算 \\
2004 & Gourdon & $10^{13}$ & 计算 \\
2020 & Platt \& Trudgian & $3\times 10^{12}$ & 严格验证 \\
\bottomrule
\end{tabular}
\end{table}

其中 Platt--Trudgian (2021) 首次在\textbf{可检验的数学证明}意义上确认：
前 $3\times 10^{12}$ 个非平凡零点均在临界线上（区间算术与围道积分给出每个零点的严格包围）。
Gourdon (2004) 的 $10^{13}$ 属于大规模\textbf{数值计算}（未附带同等意义的逐点严格证明）；
二者数量不可直接比「谁更强」——一个侧重高度与通量，一个侧重可验证的逻辑强度。
Turing 方法等经典技术为早期计算提供了理论框架；现代实现则广泛使用
RS 型求值、并行与（在严格路线上）球算术，与上文工具分层一致。

% ============================================================
\section{数值实验}\label{sec:numerical}
% ============================================================

\subsection{零点计算示例}

利用第~\ref{sec:zeros}~节方法，计算前十个零点虚部如表~\ref{tab:zeros}。
对于数值逼近实验中采用的共轭非平凡零点个数分别取 $K = 20,\, 60,\, 200,\, 500$ 的情形，
其对应的最高非平凡零点虚部高度分别为：
$\gamma_{20} \approx 77.15$,
$\gamma_{60} \approx 163.03$,
$\gamma_{200} \approx 396.38$,
$\gamma_{500} \approx 811.18$。
在普通个人计算机上计算这些非平凡零点（以计算至最高的 $\gamma_{500}$ 为例）仅需数秒，
利用割线法可以快速精确到 $10^{-15}$ 级别。

\subsection{不同零点个数的逼近效果对比}\label{subsec:comparison_plots}

\definecolor{blueframe}{RGB}{30, 90, 160}
\definecolor{bluebg}{RGB}{235, 243, 255}
\definecolor{greenbg}{RGB}{235, 250, 240}
\definecolor{greenframe}{RGB}{30, 140, 80}
\definecolor{amberbg}{RGB}{255, 248, 225}
\definecolor{amberframe}{RGB}{180, 120, 0}

\vspace{0.1cm}

\paragraph{说明与计算参数}
$\pi_0$ 逼近（实用记 $\pi_K$）的意义在于对 $\pi(x)$ 的\textbf{整体波形拟合}
（平台贴合、跳跃被振荡重构），而\textbf{不是}在个别素数跳动点上与 $\pi(p)$ 比“点值有多准”。
跳动点处有限零点和收敛于中值 $\pi(p)-\frac12$，附近可有起伏——属有限项逼近阶梯的通性；
\textbf{不宜}与周期 Fourier 级数的经典 Gibbs 过冲（约 $9\%$、Dirichlet 核）混称为同一定理
（见第~\ref{ch:fourier_vs_riemann}~章）。评价逼近质量应以区间上的曲线对照为主
（图~\ref{fig:pi_comparison_merged}）。

计算采用第~\ref{sec:explicit}~节的 Möbius--$\Ei$ 截断~\eqref{eq:pi_explicit_full}，
并固定本图参数如下：
\begin{itemize}
  \item $N_{\mu}=\lfloor\log_2 x\rfloor+1$；
  \item 平凡零点截断 $M = 10$；
  \item 非平凡零点共轭对数 $K = 20,\, 60,\, 200,\, 500$
        （中档取 $60$ 以区别于第~\ref{ch:fourier_vs_riemann}~章 $N=100$ 的对照图）；
  \item 绘图区间 $x \in [10, 40]$（含 $11,13,17,19,23,29,31,37$ 等跳跃）；
  \item 蓝线数据文件 \texttt{riemann\_data\_*.dat} 中的纵坐标相对 $\pi_K$ 常有固定平移
        （作图时取 $y+6$），仅为与右连续阶梯 $\pi(x)$ 同框对齐，\textbf{不}改变波形与相对误差。
\end{itemize}
逼近过程的动态说明、截图与作者提供的视频及源程序，
见第~\ref{ch:visualization}~章第~\ref{sec:explicit_vis}~节
（$\pi_0$ 逼近：从零点到素数阶梯）、书后附注，
及 GitHub 仓库 \url{https://github.com/lichengtong123/lichengtong-RH}。

\begin{figure}[htbp]
\centering

% 行 1
\begin{minipage}{0.49\linewidth}
\centering
\begin{tikzpicture}
\begin{axis}[
    width=7.2cm, height=5.2cm,
    xmin=10, xmax=40,
    ymin=3.5, ymax=13.5,
    xlabel={$x$},
    ylabel={},
    title={\small\bf (a) $K=20$ 个零点},
    xtick={11,13,17,19,23,29,31,37,40},
    xticklabels={$11$,$13$,$17$,$19$,$23$,$29$,$31$,$37$,$40$},
    ytick={4,5,6,7,8,9,10,11,12},
    yticklabels={$4$,$5$,$6$,$7$,$8$,$9$,$10$,$11$,$12$},
    grid=both,
    grid style={line width=0.2pt, draw=gray!20},
    major grid style={line width=0.3pt, draw=gray!35},
    axis lines=left,
    tick align=outside,
    legend pos=north west,
    legend style={font=\tiny, fill=white, fill opacity=0.85, draw opacity=0.5},
    tick label style={font=\tiny},
    label style={font=\tiny},
    clip=true,
]
\addplot[red!80!black, very thick, solid, const plot]
    coordinates {(10,4)(11,5)(13,6)(17,7)(19,8)(23,9)(29,10)(31,11)(37,12)(40,12)};
\addlegendentry{阶梯函数}

\addplot[blue!70!green!80!black, thick, line width=1.0pt,
    restrict y to domain=3:16]
    table[x index=0, y expr=\thisrowno{1}+6, col sep=space]
    {riemann_data_20.dat};
\addlegendentry{黎曼逼近}
\draw[gray!40, very thin, solid] (axis cs:11, 3.5) -- (axis cs:11, 13.5);
\draw[gray!40, very thin, solid] (axis cs:13, 3.5) -- (axis cs:13, 13.5);
\draw[gray!40, very thin, solid] (axis cs:17, 3.5) -- (axis cs:17, 13.5);
\draw[gray!40, very thin, solid] (axis cs:19, 3.5) -- (axis cs:19, 13.5);
\draw[gray!40, very thin, solid] (axis cs:23, 3.5) -- (axis cs:23, 13.5);
\draw[gray!40, very thin, solid] (axis cs:29, 3.5) -- (axis cs:29, 13.5);
\draw[gray!40, very thin, solid] (axis cs:31, 3.5) -- (axis cs:31, 13.5);
\draw[gray!40, very thin, solid] (axis cs:37, 3.5) -- (axis cs:37, 13.5);
\end{axis}
\end{tikzpicture}
\end{minipage}
\hfill
\begin{minipage}{0.49\linewidth}
\centering
\begin{tikzpicture}
\begin{axis}[
    width=7.2cm, height=5.2cm,
    xmin=10, xmax=40,
    ymin=3.5, ymax=13.5,
    xlabel={$x$},
    ylabel={},
    title={\small\bf (b) $K=60$ 个零点},
    xtick={11,13,17,19,23,29,31,37,40},
    xticklabels={$11$,$13$,$17$,$19$,$23$,$29$,$31$,$37$,$40$},
    ytick={4,5,6,7,8,9,10,11,12},
    yticklabels={$4$,$5$,$6$,$7$,$8$,$9$,$10$,$11$,$12$},
    grid=both,
    grid style={line width=0.2pt, draw=gray!20},
    major grid style={line width=0.3pt, draw=gray!35},
    axis lines=left,
    tick align=outside,
    legend pos=north west,
    legend style={font=\tiny, fill=white, fill opacity=0.85, draw opacity=0.5},
    tick label style={font=\tiny},
    label style={font=\tiny},
    clip=true,
]
\addplot[red!80!black, very thick, solid, const plot]
    coordinates {(10,4)(11,5)(13,6)(17,7)(19,8)(23,9)(29,10)(31,11)(37,12)(40,12)};
\addlegendentry{阶梯函数}

\addplot[blue!70!green!80!black, thick, line width=1.0pt,
    restrict y to domain=3:16]
    table[x index=0, y expr=\thisrowno{1}+6, col sep=space]
    {riemann_data_60.dat};
\addlegendentry{黎曼逼近}
\draw[gray!40, very thin, solid] (axis cs:11, 3.5) -- (axis cs:11, 13.5);
\draw[gray!40, very thin, solid] (axis cs:13, 3.5) -- (axis cs:13, 13.5);
\draw[gray!40, very thin, solid] (axis cs:17, 3.5) -- (axis cs:17, 13.5);
\draw[gray!40, very thin, solid] (axis cs:19, 3.5) -- (axis cs:19, 13.5);
\draw[gray!40, very thin, solid] (axis cs:23, 3.5) -- (axis cs:23, 13.5);
\draw[gray!40, very thin, solid] (axis cs:29, 3.5) -- (axis cs:29, 13.5);
\draw[gray!40, very thin, solid] (axis cs:31, 3.5) -- (axis cs:31, 13.5);
\draw[gray!40, very thin, solid] (axis cs:37, 3.5) -- (axis cs:37, 13.5);
\end{axis}
\end{tikzpicture}
\end{minipage}

\vspace{0.2cm}

% 行 2
\begin{minipage}{0.49\linewidth}
\centering
\begin{tikzpicture}
\begin{axis}[
    width=7.2cm, height=5.2cm,
    xmin=10, xmax=40,
    ymin=3.5, ymax=13.5,
    xlabel={$x$},
    ylabel={},
    title={\small\bf (c) $K=200$ 个零点},
    xtick={11,13,17,19,23,29,31,37,40},
    xticklabels={$11$,$13$,$17$,$19$,$23$,$29$,$31$,$37$,$40$},
    ytick={4,5,6,7,8,9,10,11,12},
    yticklabels={$4$,$5$,$6$,$7$,$8$,$9$,$10$,$11$,$12$},
    grid=both,
    grid style={line width=0.2pt, draw=gray!20},
    major grid style={line width=0.3pt, draw=gray!35},
    axis lines=left,
    tick align=outside,
    legend pos=north west,
    legend style={font=\tiny, fill=white, fill opacity=0.85, draw opacity=0.5},
    tick label style={font=\tiny},
    label style={font=\tiny},
    clip=true,
]
\addplot[red!80!black, very thick, solid, const plot]
    coordinates {(10,4)(11,5)(13,6)(17,7)(19,8)(23,9)(29,10)(31,11)(37,12)(40,12)};
\addlegendentry{阶梯函数}

\addplot[blue!70!green!80!black, thick, line width=1.0pt,
    restrict y to domain=3:16]
    table[x index=0, y expr=\thisrowno{1}+6, col sep=space]
    {riemann_data_200.dat};
\addlegendentry{黎曼逼近}
\draw[gray!40, very thin, solid] (axis cs:11, 3.5) -- (axis cs:11, 13.5);
\draw[gray!40, very thin, solid] (axis cs:13, 3.5) -- (axis cs:13, 13.5);
\draw[gray!40, very thin, solid] (axis cs:17, 3.5) -- (axis cs:17, 13.5);
\draw[gray!40, very thin, solid] (axis cs:19, 3.5) -- (axis cs:19, 13.5);
\draw[gray!40, very thin, solid] (axis cs:23, 3.5) -- (axis cs:23, 13.5);
\draw[gray!40, very thin, solid] (axis cs:29, 3.5) -- (axis cs:29, 13.5);
\draw[gray!40, very thin, solid] (axis cs:31, 3.5) -- (axis cs:31, 13.5);
\draw[gray!40, very thin, solid] (axis cs:37, 3.5) -- (axis cs:37, 13.5);
\end{axis}
\end{tikzpicture}
\end{minipage}
\hfill
\begin{minipage}{0.49\linewidth}
\centering
\begin{tikzpicture}
\begin{axis}[
    width=7.2cm, height=5.2cm,
    xmin=10, xmax=40,
    ymin=3.5, ymax=13.5,
    xlabel={$x$},
    ylabel={},
    title={\small\bf (d) $K=500$ 个零点},
    xtick={11,13,17,19,23,29,31,37,40},
    xticklabels={$11$,$13$,$17$,$19$,$23$,$29$,$31$,$37$,$40$},
    ytick={4,5,6,7,8,9,10,11,12},
    yticklabels={$4$,$5$,$6$,$7$,$8$,$9$,$10$,$11$,$12$},
    grid=both,
    grid style={line width=0.2pt, draw=gray!20},
    major grid style={line width=0.3pt, draw=gray!35},
    axis lines=left,
    tick align=outside,
    legend pos=north west,
    legend style={font=\tiny, fill=white, fill opacity=0.85, draw opacity=0.5},
    tick label style={font=\tiny},
    label style={font=\tiny},
    clip=true,
]
\addplot[red!80!black, very thick, solid, const plot]
    coordinates {(10,4)(11,5)(13,6)(17,7)(19,8)(23,9)(29,10)(31,11)(37,12)(40,12)};
\addlegendentry{阶梯函数}

\addplot[blue!70!green!80!black, thick, line width=1.0pt,
    restrict y to domain=3:16]
    table[x index=0, y expr=\thisrowno{1}+6, col sep=space]
    {riemann_data_500.dat};
\addlegendentry{黎曼逼近}
\draw[gray!40, very thin, solid] (axis cs:11, 3.5) -- (axis cs:11, 13.5);
\draw[gray!40, very thin, solid] (axis cs:13, 3.5) -- (axis cs:13, 13.5);
\draw[gray!40, very thin, solid] (axis cs:17, 3.5) -- (axis cs:17, 13.5);
\draw[gray!40, very thin, solid] (axis cs:19, 3.5) -- (axis cs:19, 13.5);
\draw[gray!40, very thin, solid] (axis cs:23, 3.5) -- (axis cs:23, 13.5);
\draw[gray!40, very thin, solid] (axis cs:29, 3.5) -- (axis cs:29, 13.5);
\draw[gray!40, very thin, solid] (axis cs:31, 3.5) -- (axis cs:31, 13.5);
\draw[gray!40, very thin, solid] (axis cs:37, 3.5) -- (axis cs:37, 13.5);
\end{axis}
\end{tikzpicture}
\end{minipage}

\caption{Möbius--$\Ei$ 截断 $\pi_K$ 在不同非平凡零点对数 $K$ 下与右连续 $\pi(x)$ 的对比
（$K=20,\,60,\,200,\,500$，$M=10$，$x\in[10,40]$）。
红线为阶梯 $\pi$；蓝线为截断逼近（数据纵轴经固定平移以便同框，见正文参数说明）。
跳跃点处理论目标为中值，非右连续 $\pi(p)$。}
\label{fig:pi_comparison_merged}
\end{figure}

\vspace{0.3cm}

\paragraph{平凡零点项与 $x\to 1^+$ 的正则化}
对应于 $\zeta(s)$ 在负偶数 $s=-2m$（$m\ge 1$）处的零点，其贡献可写成
\[
\sum_{m=1}^{\infty}\li(x^{-2m})
=
-\int_x^{\infty}\frac{\mathrm{d}t}{t(t^2-1)\log t}.
\]
当 $x\to 1^+$ 时主项 $\li(x)\to-\infty$，而该积分趋于 $+\infty$，两者在完整显式公式中相互抵消，
使 $J_0$（从而 $\pi_0$）在 $x\downarrow 1$ 处保持有限；这是公式的结构性质，而非额外假设。
量级上该项为 $O(x^{-2}(\log x)^{-1})$：在 $x=10$ 约 $-1.84\times 10^{-3}$，在 $x=16$ 约 $-6.10\times 10^{-4}$；
对图~\ref{fig:pi_comparison_merged} 的区间 $x\in[10,40]$ 影响已小于线宽。截断取 $M=10$ 即可。

\paragraph{非平凡零点与 $\pi_0$ 逼近}
非平凡零点 $\rho=\frac12+i\gamma$ 对应振荡项 $-\sum_{\rho}\li(x^{\rho})$
（$|\rho|$ 增大时振幅递减；共轭对合并后为实振荡，有限和条件收敛）。
随所用零点个数 $K$ 增加，$\pi_K$ 对 $\pi(x)$ 的逼近拟合表现为：
\begin{itemize}
  \item \textbf{跳跃附近}：在素数 $p$ 附近，各频率项的相位趋于一致叠加，使 $\pi_K$ 出现与阶梯相符的陡升
        （图中如 $x=17,19,23,\ldots$）；
  \item \textbf{平台段}：在相邻素数之间，不同 $\gamma_k$ 的振荡相互抵消，使 $\pi_K$ 贴近水平段；
  \item \textbf{有限 $K$ 的局部波纹}：跳跃点两侧可出现过冲与振荡；
        $K$ 增大时波纹变密并向间断点收缩。此为截断级数逼近阶梯的通性；
        截断核与 Fourier 的 Dirichlet 核不同，过冲比例\textbf{不必}等于经典 Gibbs 常数约 $9\%$
        （第~\ref{ch:fourier_vs_riemann}~章）。
\end{itemize}

\begin{center}
\fbox{\parbox{0.95\textwidth}{%
\centering
\vspace{0.35em}
\textbf{第~\ref{ch:riemann_numerical_approximation}~章：重要定理与核心公式汇总}\\[0.6em]
\setlength{\tabcolsep}{5pt}%
\renewcommand{\arraystretch}{1.45}%
\begin{tabular}{@{}>{\raggedright\arraybackslash}p{3.4cm}
                >{\raggedright\arraybackslash}p{10.6cm}@{}}
\toprule
\textbf{名称} & \textbf{要点（完整式见正文编号）} \\
\midrule
Hardy $Z$ 函数
  & $Z(t)=e^{i\vartheta(t)}\,\zeta\!\bigl(\tfrac12+it\bigr)\in\mathbb{R}$，
    $|Z(t)|=\bigl|\zeta(\tfrac12+it)\bigr|$；
    极坐标形式见~\eqref{eq:hardy_polar}；实值性见第~\ref{sec:hardy_RS}~节 \\[0.55em]
Riemann--Siegel 公式
  & $Z(t)=2\displaystyle\sum_{n=1}^{N_{\mathrm{RS}}}
    \dfrac{\cos\bigl(\vartheta(t)-t\log n\bigr)}{\sqrt{n}}+R(t)$，
    $N_{\mathrm{RS}}=\bigl\lfloor\sqrt{t/(2\pi)}\bigr\rfloor$，
    $R(t)=O(t^{-1/4})$；完整陈述见~\eqref{eq:RS} \\[0.55em]
零点计数 $N(T)$
  & $N(T)=\dfrac{T}{2\pi}\log\dfrac{T}{2\pi e}+O(\log T)$
    （本章备忘~\eqref{eq:NT_chap15}；推导见第~\ref{ch:zero_distribution}~章） \\[0.55em]
$\pi_K$ 显式逼近
  & Möbius--$\Ei$ 三重截断（主项 $+$ 非平凡零点 $+$ 平凡零点）；
    \textbf{完整公式见~\eqref{eq:pi_explicit_full}}，
    参数 $N_{\mu},K,M$ 见第~\ref{sec:explicit}~节 \\[0.55em]
截断误差（RH，阶）
  & $\bigl|\psi_0-\psi_T\bigr|
    =O\!\bigl(x^{1/2}\log^2(xT)/T\bigr)$，
    $T\sim 2\pi K/\log K$；见第~\ref{sec:explicit}~节 \\[0.55em]
间断点收敛
  & 素数 $p$ 处 $\pi_K(p)\to\pi(p)-\tfrac12$（中值）；
    整体逼近看图~\ref{fig:pi_comparison_merged} \\
\bottomrule
\end{tabular}\\[0.35em]
}}
\end{center}

% ============================================================
\section*{本章小结与后续展望}\label{sec:conclusion}
\addcontentsline{toc}{section}{本章小结与后续展望}
% ============================================================

本章完成两条数值链路（公式见上文汇总框）：
\begin{enumerate}
  \item \textbf{求零点}：Hardy $Z$ 与 Riemann--Siegel 公式 $\to$ Gram 定位与精化
        （第~\ref{sec:hardy_RS}~、\ref{sec:zeros}~节）；
  \item \textbf{用零点}：$\pi_0$ 逼近（实用记 $\pi_K$，Möbius--$\Ei$ 有限项）$\to$
        图~\ref{fig:pi_comparison_merged} 上对 $\pi(x)$ 的\textbf{整体}波形拟合
        （第~\ref{sec:explicit}~、\ref{sec:numerical}~节）。
\end{enumerate}
程序库与数据源见第~\ref{sec:libraries}~节；逼近误差与 RH 下振幅阶见第~\ref{sec:explicit}~节及第~\ref{ch:riemann_explicit_formula}~章，此处不重复。

下一章（第~\ref{ch:visualization}~章）将同一套对象转为几何示意：
Hardy 轨迹、共形图像，以及 $\pi_0$ 逼近的动态过程——
对本章数值曲线作几何解读。
